\section{Advanced Interval Estimation and Sample Size Determination}

In previous sections we intuitively introduced the concept of a confidence interval and evaluated basic interval estimations for the population mean. However, in real statistical analysis and data science projects, we rarely limit ourselves to studying a single population or assuming known population parameters. We frequently need to compare two treatments (such as in A/B testing for websites), evaluate proportions or conversion rates, analyze data variability or dispersion, and, critically before any study begins, accurately calculate the sample size required to achieve a desired statistical precision or power.

In this section we formally, thoroughly, and rigorously develop the theory of interval estimation for differences in means, standard errors, proportions, variances, and sample size design.

\subsection{Estimation of the Mean and Difference Between Two Means}

When two populations are studied, primary interest often centers on estimating the difference between their means, $\mu_1 - \mu_2$, using independent or paired samples drawn from each population. The natural, minimum-variance point estimator for this difference is the difference between sample means: $\bar{X}_1 - \bar{X}_2$.

The construction of the confidence interval for $\mu_1 - \mu_2$ depends critically on the sampling design and the assumptions or knowledge regarding the population variances ($\sigma_1^2, \sigma_2^2$).

\begin{teorema}[Case 1: Independent Samples with Known Variances]
	Let $X_{11}, X_{12}, \ldots, X_{1n_1}$ and $X_{21}, X_{22}, \ldots, X_{2n_2}$ be two independent random samples of sizes $n_1$ and $n_2$ drawn from populations with means $\mu_1, \mu_2$ and known variances $\sigma_1^2, \sigma_2^2$. If the populations are normal or the sample sizes are large ($n_1, n_2 \geq 30$ by the Central Limit Theorem), a $(1-\alpha)100\%$ confidence interval for the difference of means is:
	\begin{align}
		(\bar{X}_1 - \bar{X}_2) - Z_{\alpha/2} \sqrt{\frac{\sigma_1^2}{n_1} + \frac{\sigma_2^2}{n_2}} \leq \mu_1 - \mu_2 \leq (\bar{X}_1 - \bar{X}_2) + Z_{\alpha/2} \sqrt{\frac{\sigma_1^2}{n_1} + \frac{\sigma_2^2}{n_2}},
	\end{align}
	where $Z_{\alpha/2}$ is the critical value from the standard normal distribution leaving an upper tail area of $\alpha/2$.
\end{teorema}

\begin{teorema}[Case 2: Independent Samples with Unknown but Equal Variances]
	If the population variances are unknown but homoscedasticity can be assumed ($\sigma_1^2 = \sigma_2^2 = \sigma^2$), we pool the information from both samples to obtain a degrees-of-freedom-weighted estimate of the common variance, denoted by the pooled variance $S_p^2$:
	\begin{align}
		S_p^2 = \frac{(n_1 - 1)S_1^2 + (n_2 - 1)S_2^2}{n_1 + n_2 - 2},
		\label{eq:varianza_combinada}
	\end{align}
	where $S_1^2$ and $S_2^2$ are the respective unbiased sample variances. 
	
	The standardized statistic follows a Student's $t$-distribution with $\nu = n_1 + n_2 - 2$ degrees of freedom. The $(1-\alpha)100\%$ confidence interval for $\mu_1 - \mu_2$ is:
	\begin{align}
		(\bar{X}_1 - \bar{X}_2) \pm t_{\alpha/2, \, n_1+n_2-2} \, S_p \sqrt{\frac{1}{n_1} + \frac{1}{n_2}}.
	\end{align}
\end{teorema}

\begin{teorema}[Case 3: Independent Samples with Unknown and Unequal Variances (Welch)]
	When population variances are unknown and it is not reasonable to assume they are equal ($\sigma_1^2 \neq \sigma_2^2$), we cannot use the pooled variance. In this scenario (known as the Behrens-Fisher problem), we use the individual sample variances and approximate the distribution of the test statistic using Student's $t$-distribution with Welch-Satterthwaite effective degrees of freedom:
	\begin{align}
		\nu = \frac{\left( \frac{S_1^2}{n_1} + \frac{S_2^2}{n_2} \right)^2}{\frac{\left(S_1^2 / n_1\right)^2}{n_1 - 1} + \frac{\left(S_2^2 / n_2\right)^2}{n_2 - 1}}.
		\label{eq:df_welch}
	\end{align}
	If $\nu$ is not an integer, it is typically rounded down to the nearest integer to maintain a conservative attitude regarding interval coverage. The $(1-\alpha)100\%$ interval is given by:
	\begin{align}
		(\bar{X}_1 - \bar{X}_2) \pm t_{\alpha/2, \, \nu} \sqrt{\frac{S_1^2}{n_1} + \frac{S_2^2}{n_2}}.
	\end{align}
\end{teorema}

\begin{definicion}[Case 4: Paired or Dependent Sampling]
	When two samples are pairwise correlated (for example, measurements from the same patient or system before and after a treatment), we cannot treat them as independent. Instead, we reduce the two-variable problem to a single-variable problem by computing individual paired differences:
	\begin{align}
		D_i = X_{1i} - X_{2i}, \quad \text{for } i = 1, 2, \ldots, n.
	\end{align}
	
	We compute the sample mean of the differences $\bar{D} = \frac{1}{n}\sum D_i$ and the sample standard deviation of the differences $S_D = \sqrt{\frac{1}{n-1}\sum (D_i - \bar{D})^2}$. The $(1-\alpha)100\%$ confidence interval for the population mean difference $\mu_D = \mu_1 - \mu_2$ is:
	\begin{align}
		\bar{D} - t_{\alpha/2, \, n-1} \frac{S_D}{\sqrt{n}} \leq \mu_D \leq \bar{D} + t_{\alpha/2, \, n-1} \frac{S_D}{\sqrt{n}}.
	\end{align}
\end{definicion}

\subsection{Standard Errors of Main Estimators}

The \emph{Standard Error} ($SE$) of any point estimator $\hat{\theta}$ is the standard deviation of the sampling distribution of that estimator: $SE(\hat{\theta}) = \sqrt{\text{Var}(\hat{\theta})}$. It forms the fundamental building block for the margin of error in every confidence interval.

For systematic reference in data science, Table \ref{tab:errores_estandar} summarizes the theoretical standard errors and their empirical sample estimators for the most relevant parameters in statistical inference.

\begin{table}[h!]
	\centering
	\begin{tabular}{lccc}
		\hline
		\textbf{Parameter ($\theta$)} & \textbf{Estimator ($\hat{\theta}$)} & \textbf{Theoretical Standard Error} & \textbf{Standard Error Estimate ($\widehat{SE}$)} \\ \hline
		One-pop. mean ($\mu$) & $\bar{X}$ & $\frac{\sigma}{\sqrt{n}}$ & $\frac{s}{\sqrt{n}}$ \\ [1.5ex]
		Indep. means diff. ($\mu_1 - \mu_2$) & $\bar{X}_1 - \bar{X}_2$ & $\sqrt{\frac{\sigma_1^2}{n_1} + \frac{\sigma_2^2}{n_2}}$ & $\sqrt{\frac{S_1^2}{n_1} + \frac{S_2^2}{n_2}}$ \text{ or } $S_p \sqrt{\frac{1}{n_1} + \frac{1}{n_2}}$ \\ [1.5ex]
		Paired means diff. ($\mu_D$) & $\bar{D}$ & $\frac{\sigma_D}{\sqrt{n}}$ & $\frac{S_D}{\sqrt{n}}$ \\ [1.5ex]
		Pop. proportion ($p$) & $\hat{p} = \frac{X}{n}$ & $\sqrt{\frac{p(1-p)}{n}}$ & $\sqrt{\frac{\hat{p}(1-\hat{p})}{n}}$ \\ [1.5ex]
		Proportions diff. ($p_1 - p_2$) & $\hat{p}_1 - \hat{p}_2$ & $\sqrt{\frac{p_1(1-p_1)}{n_1} + \frac{p_2(1-p_2)}{n_2}}$ & $\sqrt{\frac{\hat{p}_1(1-\hat{p}_1)}{n_1} + \frac{\hat{p}_2(1-\hat{p}_2)}{n_2}}$ \\ \hline
	\end{tabular}
	\caption{Summary of standard errors for main estimators in inferential statistics.}
	\label{tab:errores_estandar}
\end{table}

\subsection{Estimation of a Proportion and Difference Between Two Proportions}

The study of population proportions is vital in analyzing categorical binomial or Bernoulli variables, forming the foundation of classification metrics, click-through rates in e-commerce, and preference percentages.

\begin{teorema}[Confidence Interval for a Proportion (Wald Normal Approximation)]
	Let $X$ be the number of successes in a random sample of size $n$ from a Bernoulli population with true proportion $p$. The point estimator is the sample proportion $\hat{p} = X/n$. By the Central Limit Theorem, if the large-sample conditions hold ($n\hat{p} \geq 5$ and $n(1-\hat{p}) \geq 5$, or more conservatively $\geq 10$), the sampling distribution of $\hat{p}$ is approximately normal.
	
	The $(1-\alpha)100\%$ confidence interval is given by:
	\begin{align}
		\hat{p} - Z_{\alpha/2} \sqrt{\frac{\hat{p}(1-\hat{p})}{n}} \leq p \leq \hat{p} + Z_{\alpha/2} \sqrt{\frac{\hat{p}(1-\hat{p})}{n}}.
	\end{align}
\end{teorema}

\begin{observacion}
	For small or moderate sample sizes, the Wald interval can exhibit actual coverage rates significantly lower than the nominal $(1-\alpha)$ level due to the discrete skewness of the binomial distribution. In such scenarios, data scientists prefer the \emph{Wilson (or Score) Interval} or the Agresti-Coull approximation (adding 2 pseudo-successes and 2 pseudo-failures to the sample, using $\tilde{n} = n + 4$ and $\tilde{p} = (X+2)/(n+4)$).
\end{observacion}

\begin{teorema}[Wilson (Score) Interval for a Proportion]
	Unlike the Wald interval, which centers the interval at $\hat{p}$ and approximates the standard error by $\sqrt{\hat p(1-\hat p)/n}$, the Wilson interval is obtained by directly inverting the score test $Z = (\hat{p}-p)/\sqrt{p(1-p)/n}$ without substituting $p$ by $\hat{p}$ in the denominator. This yields an interval centered at a point different from $\hat p$, with asymmetric width:
	\begin{align}
		\frac{1}{1+\frac{Z_{\alpha/2}^2}{n}}\left[\hat{p} + \frac{Z_{\alpha/2}^2}{2n} \;\pm\; Z_{\alpha/2}\sqrt{\frac{\hat{p}(1-\hat{p})}{n} + \frac{Z_{\alpha/2}^2}{4n^2}}\right].
	\end{align}
\end{teorema}

\begin{observacion}[Why the Wilson interval is preferable]
	The Wilson interval maintains actual coverage much closer to the nominal $(1-\alpha)$ level even for small $n$ or $\hat p$ near $0$ or $1$ (where Wald can collapse to an interval with limits outside $[0,1]$ or near-zero width). Modern statistical software (\texttt{statsmodels}, \texttt{R}) uses it as the default or recommended option over Wald.
\end{observacion}

\begin{teorema}[Confidence Interval for the Difference Between Two Proportions]
	To compare success rates $p_1$ and $p_2$ of two independent populations, we draw samples of sizes $n_1$ and $n_2$ and compute sample proportions $\hat{p}_1 = X_1/n_1$ and $\hat{p}_2 = X_2/n_2$. If both samples are sufficiently large ($n_i\hat{p}_i \geq 5$ and $n_i(1-\hat{p}_i) \geq 5$ for $i=1,2$), the $(1-\alpha)100\%$ confidence interval for $p_1 - p_2$ is:
	\begin{align}
		(\hat{p}_1 - \hat{p}_2) \pm Z_{\alpha/2} \sqrt{\frac{\hat{p}_1(1-\hat{p}_1)}{n_1} + \frac{\hat{p}_2(1-\hat{p}_2)}{n_2}}.
	\end{align}
\end{teorema}

\subsection{Estimation of Variance and the Ratio of Two Variances}

In many statistical applications, process variability is just as crucial as its average. For instance, in quality control, quantitative finance (where variance measures asset risk or volatility), or verifying assumptions for ANOVA and linear regression.

\begin{teorema}[Confidence Interval for a Population Variance $\sigma^2$]
	Let $X_1, X_2, \ldots, X_n$ be a random sample of size $n$ from a normally distributed population with mean $\mu$ and variance $\sigma^2$. The sample statistic:
	\begin{align}
		\chi^2 = \frac{(n-1)S^2}{\sigma^2}
	\end{align}
	follows a chi-square ($\chi^2$) distribution with $n-1$ degrees of freedom.
	
	Therefore, a $(1-\alpha)100\%$ confidence interval for the population variance $\sigma^2$ (and taking the square root for the standard deviation $\sigma$) is obtained by solving for the quantiles of the $\chi^2$ distribution:
	\begin{align}
		\frac{(n-1)S^2}{\chi^2_{\alpha/2, \, n-1}} \leq \sigma^2 \leq \frac{(n-1)S^2}{\chi^2_{1-\alpha/2, \, n-1}},
	\end{align}
	where $\chi^2_{\alpha/2, n-1}$ and $\chi^2_{1-\alpha/2, n-1}$ are the upper and lower quantiles of the chi-square distribution leaving upper tail areas of $\alpha/2$ and $1-\alpha/2$, respectively.
\end{teorema}

\begin{teorema}[Confidence Interval for the Ratio of Two Variances $\sigma_1^2 / \sigma_2^2$]
	Let two independent random samples of sizes $n_1$ and $n_2$ be drawn from independent normal populations with variances $\sigma_1^2$ and $\sigma_2^2$. The ratio of the standardized sample variances:
	\begin{align}
		F = \frac{S_1^2 / \sigma_1^2}{S_2^2 / \sigma_2^2} = \left(\frac{S_1^2}{S_2^2}\right)\left(\frac{\sigma_2^2}{\sigma_1^2}\right)
	\end{align}
	follows Snedecor's $F$-distribution with $\nu_1 = n_1 - 1$ numerator degrees of freedom and $\nu_2 = n_2 - 1$ denominator degrees of freedom.
	
	The $(1-\alpha)100\%$ confidence interval for the population variance ratio $\frac{\sigma_1^2}{\sigma_2^2}$ is:
	\begin{align}
		\frac{S_1^2}{S_2^2} \frac{1}{F_{\alpha/2, \, n_1-1, \, n_2-1}} \leq \frac{\sigma_1^2}{\sigma_2^2} \leq \frac{S_1^2}{S_2^2} \frac{1}{F_{1-\alpha/2, \, n_1-1, \, n_2-1}},
	\end{align}
	where $F_{\alpha/2, \nu_1, \nu_2}$ is the upper quantile of the $F$-distribution. 
	
	Recalling the reciprocal property of the $F$-distribution, we have $F_{1-\alpha/2, \nu_1, \nu_2} = \frac{1}{F_{\alpha/2, \nu_2, \nu_1}}$, which simplifies calculations when statistical tables report only upper tails.
\end{teorema}

\subsection{Sample Size Determination}

One of the questions with the highest economic and computational impact in experimental data science is: \emph{How much data do we need?} Taking too small a sample may leave the study underpowered to detect real effects, whereas taking an excessively large sample wastes financial resources, time, and computational infrastructure.

To determine the appropriate sample size $n$ before collecting data, three parameters must be specified:
\begin{enumerate}
	\item The \textbf{nominal confidence level $(1-\alpha)$}, which determines the critical quantile ($Z_{\alpha/2}$ or $t_{\alpha/2}$).
	\item The \textbf{maximum tolerable margin of error ($E$)}, representing half the desired width of the confidence interval.
	\item A \textbf{prior estimate of population dispersion or variability} ($\sigma$ for means, or an expected proportion $p^*$ for proportions).
\end{enumerate}

\begin{teorema}[Sample Size for Estimating a Population Mean $\mu$]
	When estimating the mean of a normal population or under the Central Limit Theorem, the margin of error of the confidence interval is $E = Z_{\alpha/2} \frac{\sigma}{\sqrt{n}}$. Solving for $n$, we obtain the required sample size:
	\begin{align}
		n = \left( \frac{Z_{\alpha/2} \cdot \sigma}{E} \right)^2.
		\label{eq:n_media}
	\end{align}
	If the population variance $\sigma^2$ is unknown, a preliminary estimate can be obtained through: (a) a prior pilot study with a small sample ($n_0 \approx 30$) to use the sample standard deviation $s$; or (b) the empirical range rule ($\sigma \approx \frac{\text{Estimated Range}}{4}$ if the population is roughly normal). Whenever the result is not an exact integer, the computed value is always rounded \emph{up to the next integer} to guarantee that the final margin of error is less than or equal to $E$.
\end{teorema}

\begin{teorema}[Sample Size for Estimating a Population Proportion $p$]
	The margin of error for a population proportion using the normal approximation is $E = Z_{\alpha/2} \sqrt{\frac{p(1-p)}{n}}$. Solving for $n$, we obtain:
	\begin{align}
		n = \left( \frac{Z_{\alpha/2}}{E} \right)^2 p^* (1 - p^*),
		\label{eq:n_proporcion}
	\end{align}
	where $p^*$ is a prior estimate of the true proportion.
	
	\textbf{Conservative case (no prior information):} If no preliminary data on the likely value of $p$ are available, we must protect against the worst-case scenario regarding variability. The quadratic function $g(p) = p(1-p)$ reaches its absolute maximum at $p = 0.5$, where $g(0.5) = 0.5 \times 0.5 = 0.25$. Substituting this maximum into the previous equation yields the safest, most conservative formula:
	\begin{align}
		n_{\text{max}} = \frac{Z_{\alpha/2}^2}{4 E^2}.
		\label{eq:n_conservador}
	\end{align}
\end{teorema}

\begin{observacion}
	When the total study population is of finite and relatively small size (denoted by $N$), and the sample size computed by the formulas above represents more than $5\%$ of the population ($n / N > 0.05$), we must apply the \emph{Finite Population Correction Factor} (FPCF). The adjusted sample size $n_a$ is:
	\begin{align}
		n_a = \frac{n}{1 + \frac{n-1}{N}}.
	\end{align}
	This adjustment reduces the number of required observations without sacrificing the stipulated statistical precision.
\end{observacion}
